\section{Structural reductions for the canonical index}
\label{sec:reductions}

We now study the absolute index problem in dimension $n$, still assuming
$I_j$ and $L_j$ for $j<n$.  Proposition~\ref{prop:denominator} and
Proposition~\ref{prop:rc-torsion} already control a variety admitting a
nontrivial contraction with rationally connected base.  The purpose of
this section is to describe what an unbounded-index sequence would have
to look like.  The reductions adapt those of
\cite[Section \emph{Structural reductions for the canonical index}]{U4};
the proofs below explain the changes needed in arbitrary dimension.

For a normal projective variety with $K_X\sim_{\Q}0$, its
\emph{canonical index} is the least positive integer $r$ such that
$rK_X$ is an integral principal divisor.  In particular, this records
both Cartierness and linear triviality.

\begin{proposition}[The remaining index sequence]
\label{prop:index-reduction}
If the canonical indices of projective klt $n$-folds with
$K_X\sim_{\Q}0$ are unbounded, there is a sequence of projective
$\Q$-factorial terminal $n$-folds $V_i$, of canonical indices
$r_i\to\infty$, satisfying the following properties.
\begin{enumerate}
\item A smooth projective resolution of $V_i$ has irregularity zero,
and $\Bir(V_i)$ is countable.
\item Every positive-dimensional rational image of $V_i$ of dimension
less than $n$ has rationally connected smooth projective resolution.
\item Every effective Weil divisor on $V_i$ with positive Iitaka
dimension is big.
\item Let $\pi_i:Y_i\to V_i$ be the cyclic index cover, with its deck
group $G_i$.  Then $Y_i$ is terminal and $K_{Y_i}\sim0$.  Every
$G_i$-equivariant rational fibration of $Y_i$ having positive-dimensional
base and fibre has general-type smooth geometric generic fibre.
The same assertion holds for rational fibrations of $V_i$, with trivial
group action.
\end{enumerate}
Here a rational fibration has geometrically integral generic fibre, or
equivalently has a relatively algebraically closed base function field.
\end{proposition}

\subsection{Product covers and rational images}

We first dispose of the product cases.  Take the quasi-\'etale
Beauville--Bogomolov cover
\[
                         P=A\times\prod_j B_j\longrightarrow X
\]
from Lemma~\ref{lem:cover-comparison}, counting a positive-dimensional
abelian factor as one block.  Suppose that there are at least two
positive-dimensional blocks.  On the normal Galois comparison cover
$R\to X$, the same lemma supplies a subgroup of bounded index
preserving the individual factor fields.  Its transformations on the
associated normal Stein factors $R_j$ are regular, and each
$\dim R_j<n$.  Their volume characters have uniformly bounded order
by Lemma~\ref{lem:finite-character}, applied with empty boundary.
The volume form on $R$ is the product of the pulled-back factor forms.
Consequently its character on that subgroup has bounded order, and
multiplication by the least common multiple of the possible subgroup
indices kills the character of the whole Galois group.

An invariant power of the volume form descends to $X$ and generates its
divisorial pluricanonical sheaf: the cover is \'etale at generic points
of prime divisors, so descent introduces no zero or pole there.
Conversely, a generating pluriform on $X$ pulls back to an invariant
power of the volume form.  Thus the order of the volume character is
exactly the canonical index, which is bounded in the product case.

If $P=A$ is purely abelian, the Galois comparison cover $R\to A$ is
\'etale by purity and $R$ is again abelian.  The canonical character
is an eigenvalue on $H^n(R,\Z)$, of rank $\binom{2n}{n}$.  If that
eigenvalue has order $e$, its cyclotomic polynomial divides the integral
characteristic polynomial, and hence
\[
                         \varphi(e)\le\binom{2n}{n}.
\]
There are finitely many such integers $e$.  Their least common multiple
therefore bounds the canonical index in this case as well.

Only the case of a single nonabelian block can occur in an
unbounded-index sequence.  We record its rational-image property
separately, since this is the geometric reason that the relative index
theorem applies later.

\begin{lemma}[Rational images of a single block]
\label{lem:single-block-images}
Suppose that a projective klt $n$-fold $X$ with $K_X\sim_{\Q}0$
has a finite quasi-\'etale cover by one irreducible Calabi--Yau or
irreducible holomorphic symplectic block $P$.  Then a smooth resolution
of $X$ has irregularity zero, and every positive-dimensional rational
image of dimension less than $n$ has rationally connected smooth
projective resolution.
\end{lemma}

\begin{proof}
A regular one-form on a smooth projective model of $X$ pulls back
injectively to a reflexive one-form on $P$.  Indeed, the rational map
from the normal variety $P$ to the proper smooth target is defined at
every codimension-one point.  The pulled-back form is regular there
and extends reflexively.  The block has no reflexive one-forms, so
Hodge symmetry gives irregularity zero.

It remains to exclude a dominant rational map to a non-uniruled smooth
projective variety $T$ with $0<j=\dim T<n$.  By the
pseudo-effectivity criterion of BDPP~\cite{BDPP2013}, $K_T$ is
pseudo-effective.  Nonvanishing in LA gives a nonzero
$\eta\in H^0(T,mK_T)$ for some $m>0$.  Resolving the composite map
from $P$, we obtain $u:W\to P$ and $f:W\to T$, with $W$ smooth.
Pullback of differentials gives a nonzero section
\[
             s=f^*\eta\in H^0\bigl(W,(\Omega_W^j)^{\otimes m}\bigr).
\]
At the generic point it has the form $s=a\beta^{\otimes m}$, where
$a\in\C(P)^*$ and $\beta$ is a nonzero decomposable rational
$j$-form, obtained by pulling back a rational top form on $T$.

The divisorial zero locus of $s$ is exceptional over $P$.  To prove
this, let $H=u^*H_P$ for an ample Cartier divisor $H_P$ on $P$.
Canonicality and $K_P\sim0$ imply that $K_W$ is rationally linearly
equivalent to an effective $u$-exceptional divisor.  In particular,
\[
                            K_W H^{n-1}=0.
\]
Generic semipositivity of the cotangent bundle~\cite{CampanaPaun2015}
implies that every torsion-free quotient of
$\mathcal V=(\Omega_W^j)^{\otimes m}$ has nonnegative degree against
$H^{n-1}$.  Here one first uses an ample perturbation $H+\epsilon H_0$.
On a sufficiently high complete-intersection curve, the
Harder--Narasimhan factors of $\Omega_W^1$ restrict to semistable
bundles of nonnegative slope~\cite{MehtaRamanathan1982}.
Tensor products of semistable bundles in characteristic zero are
semistable~\cite{RamananRamanathan1984}, and exterior powers are
quotients of tensor powers.  This proves the assertion for
$\mathcal V$ with the perturbed polarization; passage to the limit
gives the assertion for $H$.

The saturation of the line generated by $s$ is $\OO_W(Z)$, where
$Z\ge0$ is its divisorial zero divisor.  Since
$c_1(\mathcal V)$ is a positive multiple of $K_W$, the quotient
inequality gives $ZH^{n-1}\le0$.  Every nonexceptional prime has
strictly positive $H^{n-1}$-degree.  Thus $Z$ is exceptional.

At the generic point of a prime divisor of $P$, write $\beta$ in a
regular frame and let $b$ be the minimum order of its coefficients.
The absence of zeros and poles of $s$ there gives
\[
                              \ord(a)+mb=0.
\]
All coefficients of $\Div_P(a)$ are therefore divisible by $m$.
Normalize $P$ in a field component obtained by adjoining an $m$th
root $z$ of $a$.  This is a finite quasi-\'etale cover: over a
divisorial discrete valuation, removing the uniformizer to its
$m$-divisible order leaves the equation for a root of a unit, which
is unramified in characteristic zero.  The rational form $z\beta$
has no divisorial pole, hence is a nonzero reflexive $j$-form on the
cover.  It remains decomposable.

This contradicts the defining form algebra of the block.  In the
Calabi--Yau case there are no intermediate-degree reflexive forms on
any connected quasi-\'etale cover.  In the symplectic case the only
possibilities are scalar multiples of $\sigma^q$ with $0<2q<n$,
where $\sigma$ is generically nondegenerate.  Contraction into
$\sigma^q$ is injective on the nondegenerate locus, whereas a
nonzero decomposable $j$-form with $j<n$ has a contraction kernel of
dimension $n-j$.  These forms cannot agree.

Finally, a smooth resolution of any proper rational image has a
maximal rationally connected quotient whose smooth projective base
is non-uniruled~\cite{GHS2003}.  A positive-dimensional such base would
be a forbidden rational image of $X$.  The quotient is therefore a
point, proving rational connectedness.
\end{proof}

\subsection{Terminality and moving divisors}

If a single-block quotient $X$ is not canonical, choose a smooth
resolution $W\to X$ exhibiting a negative discrepancy.  Then
$K_W\sim_{\Q}E$ for an exceptional divisor $E$ having a negative
coefficient.  The canonical class of $W$ is not pseudo-effective.
Indeed, otherwise LA gives a nonzero pluricanonical section and hence
an effective divisor rationally linearly equivalent to $E$.  Pushing
that relation down to $X$ shows that its rational function has
effective principal divisor.  Projectivity forces that function to
be constant, so the effective divisor would equal $E$, a contradiction.
BDPP now makes $W$ uniruled.  Its maximal rationally connected
quotient has dimension less than $n$; Lemma~\ref{lem:single-block-images}
forces its base to be a point.  Proposition~\ref{prop:rc-torsion},
with zero boundary and zero nef part, bounds the index of $X$.

We may consequently discard all noncanonical terms of an
unbounded-index sequence.  A canonical variety has a projective
crepant $\Q$-factorial terminalization by the extraction theorem of
BCHM~\cite{BCHM2010}.  On a fixed resolution there are only finitely many
exceptional divisors of discrepancy zero: any valuation exceptional
over that smooth resolution has positive discrepancy, and adding
the pullback of its effective discrepancy divisor preserves
positivity.  Extracting precisely those zero-discrepancy divisors,
together with a small $\Q$-factorialization when necessary, gives
the desired terminalization.  Crepant pullback and divisor pushforward
preserve the exact principal multiples of the canonical divisor.
Write the resulting terminal varieties as $V_i$.

We also need the countability conclusion, in the precise generality
proved in \cite[Lemma 8.2]{U4}: if a
terminal projective variety has torsion canonical class and
irregularity zero on a resolution, its birational group is countable.
Here is the mechanism.  The Picard group of a smooth resolution is
countable by the exponential sequence and $H^1(\OO)=0$; its
surjection onto $\Cl(V_i)$ makes the latter countable.  Every
birational self-map of $V_i$ is an isomorphism in codimension one.
On a common resolution the two pullbacks of a canonical
trivialization are proportional, and terminality identifies their
zero supports with the respective exceptional divisors, so those
supports coincide.  Group the self-maps by the Weil class of the
strict transform of a fixed very ample divisor.  Within one class
the maps differ by a projective linear stabilizer of the corresponding
embedding.  The effective identity component of that stabilizer is
trivial by Matsumura's theorem~\cite[Corollary~1]{Matsumura1963}
excluding a positive-dimensional
linear algebraic group of birational transformations on a smooth
complete variety with a nonzero pluricanonical system.  Each class
therefore contributes only finitely many maps.

Suppose that an effective Weil divisor $J$ on some $V_i$ has
$0<\kappa(V_i,J)<n$.  Since $V_i$ is $\Q$-factorial, we may choose
a small positive rational $\delta$ with $(V_i,\delta J)$ klt.
LA gives a good model of this pair and a semiample contraction
$f:V_i'\to Z$ with $0<\dim Z<n$.  The zero-boundary canonical
divisors remain crepant during this program.  More explicitly, fix
$rK_{V_i}=\Div(h)$ with compatible canonical representatives.  The
same relation pushes through each divisorial contraction and through
the codimension-one identification in each flip.  Since the models
are $\Q$-factorial, their canonical pullbacks to a common resolution
are all $r^{-1}\Div(h)$.  Thus $V_i'$ is still klt,
$K_{V_i'}\sim_{\Q}0$, and its canonical index is unchanged.

Apply Proposition~\ref{prop:denominator} to $(V_i',0)\to Z$.
Its generalized adjoint $D_Z$ is rationally linearly trivial.  The
base is a rational image of $V_i$, so its smooth resolution is
rationally connected by Lemma~\ref{lem:single-block-images}.
Proposition~\ref{prop:rc-torsion} therefore gives a uniform principal
multiple of $D_Z$ and hence of $K_{V_i'}$.  The bound depends on
$n$ alone, not on $J$ or $\delta$, because the canonical bundle
formula was applied with boundary zero.  Discarding finitely many
terms of our sequence makes this impossible.  Every effective Weil
divisor of positive Iitaka dimension on each remaining $V_i$ is big.

\begin{lemma}[General-type fibres from moving divisors]
\label{lem:general-type-fibres}
Let $V$ be projective, $\Q$-factorial and terminal, with
$K_V\sim_{\Q}0$, and suppose that every effective Weil divisor of
positive Iitaka dimension is big.  Let $\pi:Y\to V$ be its cyclic
index cover with deck group $G$.  Every intermediate
$G$-equivariant rational fibration of $Y$ has general-type smooth
geometric generic fibre.  The same is true on $V$ without a group
action.
\end{lemma}

\begin{proof}
The cover is connected by minimality of the index and is quasi-\'etale.
The finite discrepancy formula makes $Y$ terminal: an exceptional
divisorial valuation over $Y$ restricts to an exceptional valuation
over $V$, and its log discrepancy is the positive ramification index
times a number greater than one.  Also $K_Y\sim0$.

Let $\C(T)\subset\C(Y)$ be the relatively algebraically closed
base field of an equivariant fibration.  Its invariant field
$\C(T)^G$ lies in $\C(V)$, has the same transcendence degree, and
is relatively algebraically closed there.  Indeed, an element of
$\C(V)$ algebraic over $\C(T)^G$ is algebraic over $\C(T)$, hence
lies in $\C(T)$ and then in its invariant field.  This gives a
descended intermediate rational fibration $V\dashrightarrow Z$.

Resolve the maps and bases, obtaining $u:U\to Y$ with $U$ smooth,
a morphism $f:U\to T'$ to the unquotiented base, and a morphism
$g:U\to Z'$ constant on the generic fibres of $f$.  Pull back a
general very ample hyperplane $H$ on $Z'$ by the rational map from
$V$, and close up the divisor to obtain $D$ on $V$.  The properness
of $Z'$ ensures that this map is defined in codimension one.  The
hyperplane pullbacks form a moving linear system with a nonconstant
section ratio; thus $\kappa(V,D)>0$ and $D$ is big.

For $q=\pi\circ u$, one has
\begin{equation}\label{eq:exceptional-fibre-class}
                              q^*D=g^*H+E,
\end{equation}
with $E\ge0$ exceptional over $Y$.  At every nonexceptional prime
the two pullbacks agree, because its image on $V$ is a prime at
which the original rational map is defined.  After fixing the
resolution, choose $H$ not to contain the generic image of any of
the finitely many exceptional primes.  Its pullback has zero
coefficient at each such prime, whereas $q^*D$ is effective.  This
proves the effectivity assertion in
\eqref{eq:exceptional-fibre-class}.

Let $F$ be the smooth geometric generic fibre of $f$.  The class
$q^*D|_F$ is big: decompose the big class $q^*D$ into an ample
rational class and an effective rational class, and restrict to the
generic fibre, which is not contained in the effective support.
Since $g^*H|_F\sim0$, the class $E|_F$ is big.  Terminality gives
\[
                         K_U\sim_{\Q}\sum_a d_a E_a,
                         \qquad d_a>0,
\]
where the sum includes every $u$-exceptional prime.  Some positive
rational multiple of $E$ is coefficientwise dominated by this sum.
Restriction and adjunction therefore make $K_F=K_U|_F$ big.  The
proof for $V$ is the same with the finite cover omitted.
\end{proof}

The preceding arguments prove Proposition~\ref{prop:index-reduction}.
They reduce an unbounded index to a sequence with countable birational
groups but particularly strong restrictions on equivariant fibrations.
We next turn those restrictions into uniform estimates for section
orders and curve degrees.

\section{Scalar order and curve estimates}
\label{sec:scalar}

Throughout this section, $A\preceq B$ means that $B-A$ is rationally
linearly equivalent to an effective rational divisor.  A restriction
of this relation is made only to a member not contained in that
effective divisor.  For a big, nef, semiample rational Cartier divisor
$P$ on an integral projective variety $V$ and a smooth point $x\in V$,
define
\[
 \gamma(P;V,x)=\sup_{\substack{k>0\text{ sufficiently divisible}\\
                      0\ne s\in H^0(\widetilde V,k\mu^*P)}}
                        \frac{\ord_x(s)}{k},
\]
where $\mu:\widetilde V\to V$ is a smooth projective resolution
that is an isomorphism near $x$, and the same letter denotes its lift.
Proper birational pushforward and the projection formula make this
independent of the resolution.  We write $\gamma(P;V)$ for its very
general value.  Taking powers permits any further fixed divisibility
in $k$.  For an ample $P$, also put
\[
 \varepsilon(P;x)=\inf_{\substack{x\in C\subset V\\C\text{ integral curve}}}
                         \frac{P\cdot C}{\operatorname{mult}_x C}.
\]
We use its very general value, denoted $\varepsilon(P)$.

Both invariants are homogeneous under positive rational rescaling.
Moreover, if $d=\dim V$, the asymptotic section count and the number
of jet conditions at a smooth point give
\begin{equation}\label{eq:gamma-volume}
                              \gamma(P;V)\ge(P^d)^{1/d}.
\end{equation}
Indeed, $h^0(kP)=P^d k^d/d!+O(k^{d-1})$ in divisible degrees,
whereas vanishing to order at least $l$ imposes at most
$\binom{d+l-1}{d}$ conditions.

\begin{proposition}[Scalar estimates]
\label{prop:scalar}
Assume $I_j$ for $j<n$.  There are constants $C_n,c_n>0$ with the
following property.  Let $V$ be a normal projective canonical
$n$-fold with $K_V\sim_{\Q}0$, and let a finite group $G$ act on
$V$.  Suppose that every $G$-equivariant rational fibration with
geometrically integral generic fibre and positive-dimensional base
and fibre has general-type smooth geometric generic fibre.  If $P$
is an ample rational Cartier divisor with $g^*P\sim_{\Q}P$ for
every $g\in G$, then
\begin{equation}\label{eq:scalar-estimates}
       \gamma(P;V)\le C_n(P^n)^{1/n},
       \qquad \varepsilon(P)\ge c_n(P^n)^{1/n}.
\end{equation}
The constants are independent of $V$, $G$, and the denominator of $P$.
\end{proposition}

The proof has two parts.  An excessive section order produces
smaller-dimensional covering members with small polarized canonical
degree.  We pass to their Iitaka fibres and repeat until their section
orders become small.  Their chain quotient then contradicts the
general-type hypothesis.  Once the upper order bound is known, a
thickening of a small-degree fibre gives the lower curve bound.

\subsection{The chain and tracking inputs}

We state the dimension-independent inputs from U4 in the form used
here.  A marked covering family on $V$ consists of an integral
parameter space $S$, a dominant mark map $u:S\to V$, and a reduced
incidence $\mathcal I\subset S\times V$.  Its fibre $C_b$ over a
geometric generic parameter $b$ is integral of positive dimension
and contains $u(b)$; the second projection evaluates its points in
$V$.  One may replace $S$ by a dense open or a dominant extension
to specify geometric components or auxiliary data.

Fix a field of definition $k_*$ for these data.  A generic movement
from a geometric generic mark $a$ chooses a generic parameter $b$
on a geometric component of $u^{-1}(a)$, and then a generic endpoint
on $C_b$.  At each step genericity is over the field generated by
$k_*$ and the entire preceding path, enlarged to define the chosen
geometric component.  In particular, new parameters remain generic
after the earlier marks, parameters, and endpoints have been fixed.

\begin{lemma}[Generic chains and order bounds]
\label{lem:chain-input}
For a finite nonempty list of marked covering families on an integral
projective variety in characteristic zero, there is a rational
fibration $\xi:V\dashrightarrow M$ with geometrically integral
generic fibre.  If its generic leaf has dimension $h$, then
$1\le h\le\dim V$, and every generic member is contained in its
leaf.  From a geometric generic mark $a$, a path of at most $h$
generic movements has endpoint generic in that leaf over
$\overline{k_*(a)}$.
The construction is equivariant for any finite group preserving the
list.  A dominant rational map carrying upper movements to lower
movements or stationary steps carries upper leaves into lower leaves.
For a finite group quotient $V\to V'$, the images of the leaves of
an invariant list are the generic leaves of a rational fibration on
$V'$.

Let $P$ be big, nef and semiample.  Suppose each geometric generic
marked member satisfies either of the following bounds, with the
same number $B$:
\begin{enumerate}
\item its mark $u(b)$ is smooth on $C_b$ and
$\gamma(P|_{C_b};C_b,u(b))\le B$;
\item it is a curve and
$P\cdot C_b/\operatorname{mult}_{u(b)}C_b\le B$.
\end{enumerate}
Let $H$ be a smooth projective resolution of the geometric generic
leaf, and retain pullback notation for $P$.  Then
\begin{equation}\label{eq:chain-input-bounds}
                         \gamma(P;H)\le hB,
                         \qquad P^hH\le(hB)^h.
\end{equation}
\end{lemma}

This is \cite[Lemmas 5.1 and 5.4 and Corollary 5.2]{U4}.
The bounds concern the entire section spaces on a resolution of the
member, not merely restrictions of ambient sections.  This distinction
is important when we apply the lemma to successive Iitaka fibres.

\begin{lemma}[Tracking and covering]
\label{lem:tracking-input}
Let $Z$ be projective, $\Q$-factorial and klt of dimension $d\ge2$,
and let $N$ be big, nef and semiample.  Suppose that $d+1$ positive
rational numbers with consecutive spacing $\Delta>0$ lie strictly
between $(N^d)^{1/d}$ and $\gamma(N;Z)$.  Then there is a covering
family of integral subvarieties $S$ of dimension $0<e<d$ whose
smooth geometric generic resolution $S^*$ satisfies
\begin{align}
 N^eS&\le(2/\Delta)^{d-e}N^d,\label{eq:tracking-degree}\\
 K_{S^*}(N|_{S^*})^{e-1}
 &\le\bigl(K_Z+(2d/\Delta)N\bigr)|_S(N|_S)^{e-1}.
 \label{eq:tracking-canonical}
\end{align}
In addition, if a projective variety has a smooth model with
nonnegative Kodaira dimension, then a smooth projective model of
the geometric generic member of any covering family of integral
subvarieties has nonnegative Kodaira dimension.
The same statement holds with ``general type'' in place of
``nonnegative Kodaira dimension.''
\end{lemma}

The tracking assertion is
\cite[Lemma 6.3]{U4};
its proof combines the parameter-differentiation method of
Ein--K\"uchle--Lazarsfeld
\cite[Proposition~2.3]{EinKuchleLazarsfeld1995} with numerical
subadjunction.  The tracking lemma in U4 is stated for arbitrary $d$.
The covering assertion is
\cite[Lemma 6.4]{U4}.
For orientation, one cuts the parameter space so that incidence
evaluation is generically finite over the ambient variety, then
resolves.  The formula $K_I=q^*K_Z+R$ with $R\ge0$, restricted to
a generic parameter fibre, proves both assertions about canonical
dimension.  The cuts are chosen over independent generic coefficients,
so they test the original geometric generic member.

All applications of these lemmas are at geometric generic data.
Finite constructions there spread after a dominant parameter
extension.  We never specialize a previously generic parameter.
Countable fields of definition accommodate simultaneously the jet
conditions in all divisible degrees.  This also allows the geometric
generic constructions to be realized over $\C$ when applying the
model and positivity theorems.

\subsection{Small volumes on Iitaka fibres}

The next lemma is the additional scalar input needed beyond
dimension four.  The corresponding lemma of U4 treats members of
dimension at most three.  Here weak positivity is combined with
flattening so that the Iitaka base may have any dimension less than
that of the member.

\begin{lemma}[Small adjoint volume on Iitaka fibres]
\label{lem:small-fibre-volume}
Fix $1\le e<n$ and $C'>0$.  Let $(W_i,L_i)$ be a sequence of
smooth projective $e$-folds with $\kappa(W_i)\ge0$ and big, nef,
semiample rational Cartier divisors $L_i$, such that
\begin{equation}\label{eq:small-volume-hypotheses}
                  L_i^e\longrightarrow0,
                  \qquad K_{W_i}L_i^{e-1}\le C'L_i^e.
\end{equation}
Assume $I_e$.  After a subsequence and resolution of the Iitaka
fibrations, their smooth geometric generic fibres $U_i$ have fixed
positive dimension, satisfy $\kappa(U_i)=0$, and obey
\begin{equation}\label{eq:small-fibre-volume}
                     \vol(K_{U_i}+L_i|_{U_i})\longrightarrow0.
\end{equation}
The restricted polarizations are big, nef and semiample.  If the
Iitaka base is a point, $U_i$ denotes the entire smooth model.
\end{lemma}

\begin{proof}
We suppress the sequence index.  First, for every fixed positive
rational $a$,
\begin{equation}\label{eq:adjoint-volume-upper}
                    \vol(aK_W+L)\le(aC'+1)^eL^e.
\end{equation}
Choose a divided general member of a sufficiently divisible free
multiple of $L/a$, giving a klt effective boundary
$\Gamma\sim_{\Q}L/a$.  The big klt adjoint $K_W+\Gamma$ has a
good model by BCHM.  On a common resolution, let $P_a$ be $a$ times
the pullback of its semiample adjoint.  The negative-map comparison
gives
\[
                        aK_W+L\sim_{\Q}P_a+F_a,
                        \qquad F_a\ge0,
\]
where $F_a$ is exceptional over that model.  Thus
$P_a^e=\vol(aK_W+L)$.  With $L$ pulled back to this resolution,
nef mixed-intersection inequalities give
\[
 (aC'+1)L^e\ge(aK_W+L)L^{e-1}
 \ge P_aL^{e-1}\ge(P_a^e)^{1/e}(L^e)^{(e-1)/e},
\]
which proves \eqref{eq:adjoint-volume-upper}.  Replacing $W$ by a
higher smooth model and retaining the pullback of $L$ changes
neither this volume nor the intersection estimates: additional
canonical terms are effective exceptional divisors.

By $I_e$, a fixed integer $m_0$ defines every Iitaka fibration under
consideration.  Pass to a subsequence where $j=\kappa(W)$ is
constant.  The case $j=e$ is impossible for large indices.  The
resolved system $|m_0K_W|$ would have an integral free moving part
with generically finite map and positive integral top intersection;
hence $\vol(K_W)\ge m_0^{-e}$, contrary to
\eqref{eq:adjoint-volume-upper}.  If $j=0$, the assertion follows
directly from that inequality with $a=1$.

Suppose henceforth that $0<j<e$.  Resolve the Iitaka system and its
base, obtaining a morphism with connected fibres between smooth
projective varieties
\[
                             f_0:W_0\longrightarrow T_0.
\]
Write $L_0$ for the pulled-back polarization.  The hyperplane class
of the system's image pulls back to a big basepoint-free Cartier
divisor $H_0$ on $T_0$, and
\begin{equation}\label{eq:iitaka-hyperplane}
                         f_0^*H_0\preceq m_0K_{W_0},
                         \qquad H_0^j\ge1.
\end{equation}
The smooth geometric generic fibre $U$ has Kodaira dimension zero
by the Iitaka-fibration theorem.  The restriction $L_0|_U$ is big:
the moving fibres meet the generically finite locus of the
semiample map of $L_0$.  Nefness and semiampleness restrict as well.

We arrange a model on which sections over a big open of the base
can be used without losing volume on $W_0$.  Apply flattening by a
modification of $T_0$, and choose a smooth projective $T$ dominating
that modification.  Let $X$ be the normalization of the flat
principal family after this base change, and resolve $X$ by $W$
without changing its smooth generic fibre, which we continue to denote
by $U$.
The resulting maps fit into
\[
\begin{array}{ccccc}
 W&\longrightarrow&X&\longrightarrow&W_0\\
 &&\downarrow&&\downarrow f_0\\
 &&T&\longrightarrow&T_0.
\end{array}
\]
Write $f:W\to T$, $\rho:W\to W_0$, $L=\rho^*L_0$, and
$H$ for the pullback of $H_0$ to $T$.  Both horizontal maps in
the top row are proper and birational.  The morphism $X\to T$
is equidimensional: flatness gives the generic fibre dimension
before normalization, and finite normalization preserves it.
Hence a prime divisor $E$ on $W$ whose image in $T$ has codimension
at least two is exceptional over $X$.  Its image in $X$ has
dimension at most $e-2$, so its image in $W_0$ also has codimension
at least two.  Thus $E$ is $\rho$-exceptional.  This property will
allow extension of sections on $W_0$.  Relation
\eqref{eq:iitaka-hyperplane} gives $f^*H\preceq m_0K_W$ on the
higher smooth model, and $H^j=H_0^j\ge1$.

Choose a general projection of the morphism defined by $H$ to
$\mathbf P^j$.  It is a generically finite morphism: a general
$(j+1)$-dimensional linear subsystem has no common zero on the
$j$-dimensional image.  Its ramification divisor yields
$K_T+(j+1)H\succeq0$.  Therefore
\begin{equation}\label{eq:relative-adjoint-domination}
       K_{W/T}+L+2f^*H
              \preceq \bigl(1+(j+3)m_0\bigr)K_W+L.
\end{equation}

Choose, separately for this fibration, a sufficiently divisible
integer $l\ge1$ so that the $l$-Veronese of the section ring of
$K_U+L|_U$ is generated in degree one.  This follows from finite
generation of the big klt adjoint ring, after representing $L|_U$
by a divided general free member.  Enlarge $l$ to clear all
rational divisors and equivalences used below.  On $W$, similarly
represent $L$ by an effective klt boundary.  Twisted weak
positivity~\cite[Theorem~1.1]{Fujino2017} gives weak
positivity of the torsion-free direct image
\[
                    E_l=f_*\OO_W\bigl(l(K_{W/T}+L)\bigr).
\]
There is $b>0$ for which
$(\operatorname{Sym}^b E_l)^{**}\otimes\OO_T(bH)$ is generated
by global sections at the generic point.  A big $H$ suffices for
this consequence of weak positivity: choose $kH=A+F$ with $A$
ample and $F\ge0$, apply weak positivity with symmetric exponent
$k\beta$ and twist $\beta A$, and multiply by the section of
$\beta F$.

For every multiple $q$ of $b$, multiplication of those global
sections, followed by fibrewise multiplication in the adjoint
ring, gives sections over a big open of $T$ of
\[
                        ql(K_{W/T}+L)+qf^*H
\]
spanning all degree-$ql$ sections on the generic fibre.  To see
the asserted domain of definition, remove only the codimension-two
sets where the relevant torsion-free sheaves differ from their
reflexive hulls or are not locally free.  The multiplication maps
are defined at every codimension-one point.  Generation is needed
only at the generic point, so its possible failure along a divisor
does not require removing that divisor.

Choose a basis over $\C(T)$ among these sections.  Its size is
$h^0(U,ql(K_U+L|_U))$.  Multiplying it by a basis of
$H^0(T,qlH)$ gives linearly independent products over $\C$:
a relation would first contradict the generic-fibre independence
over $\C(T)$ and then the independence of the base sections.
If necessary, multiply by a fixed nonzero section of
$(ql-q)H$ to place these products in degree $ql$ of the left side
of \eqref{eq:relative-adjoint-domination}.  Multiplying further by
the effective difference in that equation, after clearing
denominators, embeds them into sections of its right side over
the same big open.

Let $T^\circ\subset T$ be the big open just obtained by removing
the codimension-two defects of the relevant torsion-free sheaves.
The resulting rational pluriforms are regular on $f^{-1}(T^\circ)$.
Every prime divisor on $W$ in its complement is $\rho$-exceptional,
by the flattening construction.  Thus the forms are regular at the
generic point of the strict transform of every prime divisor on
$W_0$.  Push them to $W_0$ with their $L_0$-twist.  They have no
divisorial pole and therefore extend, by smoothness, to sections of
\[
 H^0\!\left(W_0,
 \OO_{W_0}\!\left(ql\bigl((1+(j+3)m_0)K_{W_0}+L_0\bigr)\right)\right).
\]
Their independence is preserved by this birational identification.

Put $g=e-j$.  With this fibration and $l,b$ fixed, let $q$ tend
to infinity through sufficiently divisible multiples of $b$.
The asymptotic section counts on $T$ and $U$ give
\begin{equation}\label{eq:adjoint-volume-lower}
 \vol\bigl((1+(j+3)m_0)K_{W_0}+L_0\bigr)
       \ge \frac{e!}{g!j!}\,H^j\vol(K_U+L|_U).
\end{equation}
The left side tends to zero along the original sequence by
\eqref{eq:adjoint-volume-upper}, whereas $H^j\ge1$.
This proves \eqref{eq:small-fibre-volume}.  No uniform choice of
$l$ or $b$ is needed: the asymptotic estimate is obtained separately
for each fibration before taking the sequence limit.
\end{proof}

\subsection{The upper order estimate}

\begin{proof}[Proof of the first inequality in Proposition~\ref{prop:scalar}]
The assertion in dimension one follows directly from degree, so
assume $n\ge2$.  If no constant $C_n$ exists, rational rescaling
gives a sequence with
\[
                           P^n\longrightarrow0,
                           \qquad\gamma(P;V)>1.
\]
Take small projective $\Q$-factorializations and pull back $P$.
We describe a dimension-decreasing construction.  At a stage, after
passing to a subsequence, there is a fixed dimension $d$ and data
$(Z,N)$ such that
\begin{enumerate}
\item $Z$ is projective, $\Q$-factorial and klt; $N$ is big, nef
and semiample; and $N^d\to0$;
\item a smooth model of $Z$ has nonnegative Kodaira dimension;
\item $K_Z\preceq a_0N$ for a fixed positive rational $a_0$;
\item $\gamma(N;Z)$ has a positive lower bound.
\end{enumerate}
These conditions hold initially by canonicality and crepantness
of the small modification.  They cannot hold in dimension one,
because the order of a section is bounded by its degree.

For $d\ge2$, choose $d+1$ fixed equally spaced positive rational
levels below the order lower bound.  Eventually they lie above
$(N^d)^{1/d}$.  Lemma~\ref{lem:tracking-input} gives a covering
family with a smooth projective model $W$ of its geometric generic
member, of dimension $0<e<d$.  Set $L=N|_W$, using pullback notation.
It is big, nef and semiample, and $\kappa(W)\ge0$ by the
covering assertion.  The tracking inequalities and restriction of
$K_Z\preceq a_0N$ give
\[
                         L^e\longrightarrow0,
                         \qquad K_WL^{e-1}\le C'L^e
\]
for a fixed $C'$.  Lemma~\ref{lem:small-fibre-volume} therefore
gives a positive-dimensional Iitaka fibre $U$, possibly all of
$W$, with $\kappa(U)=0$ and
$\vol(K_U+L|_U)\to0$.

Choose an effective klt boundary $\Gamma\sim_{\Q}L|_U$ from
a divided general free member, and take a $\Q$-factorial good
model $Z'$ of $K_U+\Gamma$.  Let $N'$ be its semiample big
adjoint.  Then $(N')^{\dim Z'}\to0$ and
$K_{Z'}\preceq N'$.  On a common resolution with maps
$x$ to $U$ and $y$ to $Z'$, one also has
\begin{equation}\label{eq:polarization-comparison}
                             x^*(L|_U)\preceq y^*N'.
\end{equation}
Here is the fixed-part argument for this useful comparison.
Write
\[
                    x^*(K_U+\Gamma)=y^*N'+F,
                    \qquad F\ge0,
\]
with $F$ exceptional over $Z'$.  Choose an effective
$D\sim_{\Q}K_U$, which exists because $\kappa(U)=0$.
In sufficiently divisible degrees, add a general member of the
free system of $L|_U$ to this representative.  Its pullback must
contain the fixed part $F$.  The general free member contains
none of the finitely many components of $F$, so
$F\le x^*D$.  Subtracting gives
\eqref{eq:polarization-comparison}.

If $\gamma(N';Z')\to0$ on a subsequence, stop.  Otherwise pass
to a subsequence with a positive lower bound and repeat.  All
stage conditions persist: the smooth-model Kodaira dimension of
$Z'$ is zero because it is birational to $U$, and the dimension
strictly decreases.  The impossibility of the stage conditions
in dimension one forces termination.

Trace the resulting families back to the original $V$.  This
produces covering members $A\subset V$ of positive dimension
whose smooth geometric generic models satisfy
\begin{equation}\label{eq:small-order-zero-members}
                        \kappa(A^*)=0,
                        \qquad\gamma(P;A)\longrightarrow0.
\end{equation}
We make the section-order comparison explicit.  All intermediate
constructions take place on geometric generic members and spread
after dominant parameter extension.  Fibres sweep the preceding
member, and each birational comparison is an isomorphism on an
open met by the next generic member.  Composite members therefore
map birationally to their images.  At every stage the effective
polarization comparison restricts, since a sweeping member is
not contained in its extra divisor.  On common resolutions it
gives an injection from the complete section space of the
restricted preceding polarization into that of the next one,
in divisible degrees.  Multiplication by the effective
difference preserves order at a generic point outside its
support.  Proper birational pushforward preserves these complete
section spaces.  Iteration compares the complete spaces on a
resolution of $A$ with the final spaces whose normalized order
tends to zero.  This proves \eqref{eq:small-order-zero-members},
not merely an assertion about restrictions of ambient sections.

Mark these members at generic smooth points and include all their
$G$-translates.  Their order bounds are unchanged because
$g^*P\sim_{\Q}P$.  Apply Lemma~\ref{lem:chain-input}.  If the
chain leaf had full dimension, its order estimate would imply
$\gamma(P;V)\to0$, contradicting the chosen lower bound.  The
leaf has positive dimension, so eventually it is the generic
fibre of an intermediate equivariant rational fibration.

That geometric generic leaf is covered by the Kodaira-dimension-zero
members in \eqref{eq:small-order-zero-members}.  Indeed, every
generic marked member is contained in its chain leaf; after
extension to the geometric generic point of the quotient, the incidence
marks still dominate the leaf.  These extensions preserve the
geometric generic members and their Kodaira dimension.  The
covering assertion of Lemma~\ref{lem:tracking-input} therefore
prevents the leaf from being of general type.  This contradicts
the hypothesis of Proposition~\ref{prop:scalar} and proves the
upper order bound.
\end{proof}

\subsection{The lower curve estimate}

\begin{proof}[Proof of the second inequality in Proposition~\ref{prop:scalar}]
Rationally rescale so that $1\le P^n\le2$.  The upper estimate
now gives a uniform bound $\gamma(P;V)\le D_0$.  If no positive
lower bound for $\varepsilon(P)$ exists, choose covering marked
curve families with
\[
                   \frac{P\cdot C}{\operatorname{mult}_x C}
                         \longrightarrow0,
\]
and include their $G$-translates.  Such families are obtained from
curves through very general ambient marks by spreading their
geometric generic data; after shrinking, the generic marked
ratio satisfies the same bound.  Lemma~\ref{lem:chain-input}
gives chain leaves of a fixed dimension $h$ on a subsequence,
with degrees $P^hH\to0$.  Since $P^n\ge1$, eventually
$0<h<n$.

Resolve the chain fibration and choose a very general smooth
fibre $F$ of dimension $h$.  Write $s=n-h$ and retain pullback
notation for $P$.  Its normal bundle is trivial, because $F$
is a fibre of a smooth morphism near the chosen base point.
For its ideal $\mathcal I_F$, the filtration of
$\OO/\mathcal I_F^u$ has associated graded terms
$\operatorname{Sym}^j(\mathcal I_F/\mathcal I_F^2)$ for
$0\le j<u$, of total rank $\binom{u+s-1}{s}$.  Consequently
the rank of restriction of global degree-$l$ sections to this
thickened fibre is at most
\begin{equation}\label{eq:thick-fibre-bound}
                      \binom{u+s-1}{s}\,h^0(F,lP|_F).
\end{equation}
Take $u=\lceil Dl\rceil$ with $D>D_0$.  For each fixed
fibration, in large divisible degrees the leading coefficient
in \eqref{eq:thick-fibre-bound} is
\[
                    \frac{D^s}{s!h!}(P^hF)\,l^n.
\]
It tends to zero along the sequence, while the ambient section
count has leading coefficient $P^n/n!\ge1/n!$.  Thus, first
choosing a sufficiently late sequence member and then a
sufficiently large divisible $l$, there is a nonzero ambient
section vanishing in $\mathcal I_F^u$.

At any smooth point of $F$, this section has order at least $u$,
because $\mathcal I_F\subset\mathfrak m_x$.  The fibre and a
point on it can be chosen very general in the original
birational isomorphism locus and outside all the countably many
order-exceptional sets.  Birational pullback preserves the
ambient section spaces and their orders there.  The resulting
normalized order is at least $D>D_0$, contradicting the upper
bound.  This proves the second inequality and completes
Proposition~\ref{prop:scalar}.
\end{proof}
